\section{Distinct Question and Prospective Design}

This planned observational study asks whether a frozen independently evaluated measure of relational-regulation propensity predicts subsequent everyday relational actions beyond adjacent traits, context and opportunity. M03 owns scale development and initial psychometric validation; M02 owns identifiability in synthetic generators. M04 does not assume that scale scores are controller gains or that a true human controller label is available.

An ecological comprehension/technical pilot precedes independent prospective repeated measurement. Baseline scores use frozen items and scoring; repeated prompts record naturally occurring states, actions and opportunity. The proposed main endpoint is later initiation of relational action, with exact wording, schedule, sample-size rationale, metric and numeric analysis decisions pending prospective freeze. A between-person baseline propensity association is distinct from within-person state-to-action sequences. Optional partner reports require a separately specified module. No daily observation count is treated as an independent participant count.

Generalization to unseen people or couples uses a prospectively reserved group-level test set; later-time evaluation in known people is reported separately. All dyad members remain in one group split. Preprocessing, scoring and model choices are frozen before final evaluation. The study does not repeat ARC H1/H2/H3 as its primary question, induce conflict, or manipulate partner availability. Observational predictive value is not a causal-mechanism demonstration.

If M03 fails to support defensible distinct scores, confirmatory M04 pauses or prospectively revises its question around an established construct. It cannot revise a scale on final M04 outcomes and claim independent confirmation. Dynamic-parameter comparisons remain secondary and depend on uncertainty and identification limits, rather than an assumed mapping from a questionnaire to P/PI/PID.
